\pdfvariable suppressoptionalinfo 611
\providecommand{\CRevisionRound}{2}
\documentclass[8pt]{extarticle}
\usepackage[a4paper,margin=0.56in]{geometry}
\usepackage{fontspec}
\setmainfont{Latin Modern Roman}
\setsansfont{Latin Modern Sans}
\setmonofont{Latin Modern Mono}
\usepackage[english]{babel}
\babelprovide[import]{chinese-simplified}
\babelfont[chinese-simplified]{rm}[Renderer=Harfbuzz,Language=Default]{Droid Sans Fallback}
\usepackage{amsmath,amssymb,booktabs}
\newcommand{\sn}{\operatorname{sn}}
\newcommand{\cn}{\operatorname{cn}}
\newcommand{\dn}{\operatorname{dn}}
\newcommand{\sech}{\operatorname{sech}}
\title{A Complete Elliptic Atlas for the Triaxial Euler Top:\
Periods, KKS Actions, and Fixed-Circle Obstructions}
\author{Route-A structural certificate C186}
\date{}
\begin{document}
\maketitle
\begin{abstract}
For every strictly triaxial free rigid body, every nonzero angular-momentum
sphere, and every energy, we give two explicit Jacobi charts, exact minimal
periods, all axial stability rates, the intermediate-axis heteroclinic
network, and KKS action--angle coordinates.  The same atlas classifies every
fixed component of a sampled time map.  Resonant components are whole
circles, so sufficiently large iterates do not have finite isolated
fixed-point counts.  The Hamiltonian flow has a canonical unitary Koopman
owner but no intrinsic arithmetic gate, target divisor, or Hilbert--Pólya
interpretation.
\end{abstract}
\noindent\textbf{Keywords:} Euler top; Jacobi elliptic functions; action--angle
variables; Koopman group; clean fixed sets.

\begin{otherlanguage}{chinese-simplified}
\begin{center}{\large 中文摘要}\end{center}
本文对任意严格三轴自由刚体、任意非零角动量球面及全部能量层，统一给出两套
\foreignlanguage{english}{Jacobi} 椭圆图册、精确最小周期、六个主轴平衡点的稳定性、
中间轴异宿分离网以及 \foreignlanguage{english}{KKS} 作用量与角变量。固定时间映射
在共振能量上固定整条圆，因此足够高的迭代不具有有限的孤立不动点计数。该系统虽有
典范幺正 \foreignlanguage{english}{Koopman} 动力学，但没有内生素数时钟或目标除子。

\noindent 关键词：欧拉陀螺；椭圆函数；作用量与角变量；固定圆；谱障碍。
\end{otherlanguage}

\section{Frozen reduced system}
Let $0<I_1<I_2<I_3$, set
$a=I_1^{-1}>b=I_2^{-1}>c=I_3^{-1}$, and fix $|M|=G>0$.
For $m=M/G$ the Euler equation and normalized energy are
\begin{equation}
 \dot m=Gm\times\operatorname{diag}(a,b,c)m,
 \qquad e=\frac{2E}{G^2}=am_1^2+bm_2^2+cm_3^2.       \tag{1}
\end{equation}
Hence $c\le e\le a$.  Exact Jacobi representations are classical;
we use the modern computational account in~\cite{Celledoni2008} and the
explicit $4K$ convention in~\cite{Pina2015}.  Our contribution is the
all-energy theorem, consequence ledger, executable convention audit, and
strict owner boundary, not priority for the classical solution.

\section{All-energy theorem}
\paragraph{Topology and axial modes.}
The endpoint levels $e=c,a$ consist respectively of $\pm e_3$ and
$\pm e_1$.  Every regular level on either side of $b$ is two circles:
$\operatorname{sgn}m_3$ distinguishes $c<e<b$, while
$\operatorname{sgn}m_1$ distinguishes $b<e<a$.  At $e=b$, the points
$\pm e_2$ are joined by four heteroclinic branches.  The squared nonzero
tangent rates at axes $1,2,3$ are
\begin{equation}
 -G^2(a-b)(a-c),\quad G^2(a-b)(b-c),\quad
 -G^2(a-c)(b-c),                                      \tag{2}
\end{equation}
so the outer axes are elliptic and the intermediate axis is hyperbolic.

\paragraph{Low-energy chart.}
For $c<e<b$, put
\[
\begin{gathered}
 A^2=\frac{e-c}{a-c},\quad B^2=\frac{e-c}{b-c},\quad
 C^2=\frac{a-e}{a-c},\\
 k^2=\frac{(a-b)(e-c)}{(b-c)(a-e)},\qquad
\Omega^2=G^2(b-c)(a-e).
\end{gathered}
\]
For $\sigma=\pm1$ and $u=\sigma\Omega(t-t_0)$, every orbit is
\begin{equation}
 (m_1,m_2,m_3)=(A\cn u,\ B\sn u,\ \sigma C\dn u),
 \qquad T_3(e)=\frac{4K(k)}{G\sqrt{(b-c)(a-e)}}.       \tag{3}
\end{equation}

\paragraph{High-energy chart.}
For $b<e<a$, put
\[
\begin{gathered}
 A^2=\frac{e-c}{a-c},\quad B^2=\frac{a-e}{a-b},\quad
 C^2=\frac{a-e}{a-c},\\
 k^2=\frac{(b-c)(a-e)}{(a-b)(e-c)},\qquad
\Omega^2=G^2(a-b)(e-c).
\end{gathered}
\]
Then every orbit is
\begin{equation}
 (m_1,m_2,m_3)=(\sigma A\dn u,\ B\sn u,\ C\cn u),
 \quad u=\sigma\Omega(t-t_0),\quad
 T_1(e)=\frac{4K(k)}{G\sqrt{(a-b)(e-c)}}.             \tag{4}
\end{equation}
In both charts $0<k<1$.  Although $\dn$ has real period $2K$, the vector
contains $\sn$ and $\cn$ and has minimal period $4K$.

\paragraph{Separatrix and endpoint periods.}
At $e=b$, define
\[
 A_s^2=\frac{b-c}{a-c},\quad C_s^2=\frac{a-b}{a-c},
 \qquad \rho=G\sqrt{(a-b)(b-c)}.
\]
For $\varepsilon,\sigma=\pm1$ the four branches are
\begin{equation}
 m_1=\varepsilon A_s\sech u,\quad m_2=\tanh u,\quad
 m_3=\varepsilon\sigma C_s\sech u,\quad
 u=\sigma\rho(t-t_0).                                \tag{5}
\end{equation}
Thus they are heteroclinic from one intermediate-axis rotation to the
other, not periodic orbits.  Since $K(k)\to\infty$ as $k\to1$, both regular
periods diverge at $e=b$.  At the stable endpoints,
\begin{equation}
 T_3(c)=\frac{2\pi}{G\sqrt{(b-c)(a-c)}},\qquad
 T_1(a)=\frac{2\pi}{G\sqrt{(a-b)(a-c)}}.              \tag{6}
\end{equation}

\section{Proof and action--angle atlas}
Taking scalar products of $M\times I^{-1}M$ with $M$ and $I^{-1}M$
proves both quadratic integrals.  Lagrange multipliers on the sphere give
the six critical points and the stated level topology; tangent
linearization gives (2).

For (3), substitute $\cn^2=1-\sn^2$ and
$\dn^2=1-k^2\sn^2$.  The constant and $\sn^2$ coefficients in both
quadratic integrals vanish.  The rules
\[
 \sn'=\cn\dn,\qquad \cn'=-\sn\dn,\qquad
 \dn'=-k^2\sn\cn
\]
reduce the three equations (1) to the displayed $\Omega^2$; the proof of
(4) is the same coefficient calculation with the $\dn$ component
permuted.  Taking $k\to1$ and using $\sn(u,1)=\tanh u$ and
$\cn(u,1)=\dn(u,1)=\sech u$ proves (5).  Taking $k\to0$ proves (6).

Freeze $\{F,H\}=-M\cdot(\nabla F\times\nabla H)$ and
$\dot F=\{F,H\}$, so (1) has $\dot M=M\times\nabla H$.  On a positive
low component take $q=\arg(M_1+iM_2)$ and $P_3=G-M_3$; since $G$ is a
Casimir, $\{q,P_3\}=1$.  With $A_q=a\cos^2q+b\sin^2q$, its unsigned cap
action $J_3=(2\pi)^{-1}\oint P_3\,dq$ is
\begin{equation}
 J_3(e)=G\left[1-\frac1{2\pi}\int_0^{2\pi}
 \sqrt{\frac{A_q-e}{A_q-c}}\,dq\right].              \tag{7}
\end{equation}
For a positive high component take $q=\arg(M_2+iM_3)$ and
$P_1=G-M_1$, so $\{q,P_1\}=1$.  With
$B_q=b\cos^2q+c\sin^2q$, $J_1=(2\pi)^{-1}\oint P_1\,dq$ is
\begin{equation}
 J_1(e)=G\left[1-\frac1{2\pi}\int_0^{2\pi}
 \sqrt{\frac{e-B_q}{a-B_q}}\,dq\right].              \tag{8}
\end{equation}
The negative components have the same unsigned action.  Differentiating
enclosed symplectic area gives the one-degree-of-freedom identity
$|dJ/dE|=T/(2\pi)$.  Finally
$\theta=\pi u/(2K(k))\pmod{2\pi}$ satisfies
$\dot\theta=\pm2\pi/T(e)$.

\section{Sampled time, Koopman owner, and Route-A stop}
Let $\Phi_\tau$ be the time-$\tau$ map, $\tau>0$.  Its $n$th iterate fixes
all six equilibria.  A regular energy component is fixed pointwise exactly
when
\begin{equation}
 n\tau=qT(e)\qquad(q\in\mathbb Z_{\ge1}),             \tag{9}
\end{equation}
and no interior separatrix point is fixed.  The periods are continuous,
finite at the stable endpoints, and divergent at $b$.  Therefore for every
$\tau>0$, all sufficiently large iterates have at least one fixed circle.
The full-sphere isolated-cardinality Artin--Mazur coefficients are thus not
finite; selecting one energy circle would change the frozen owner.

KKS area preservation makes $U_t f=f\circ\Phi_t$ a canonical unitary
Koopman group.  Off the measure-zero critical levels it is a direct integral
over the two regimes and two components; circle Fourier mode $\ell$ carries
the phase $\exp(2\pi i\ell t/T(e))$.  This is a natural A4 coordinate.
It does not supply rational primes, prime powers, a logarithmic clock, a
target divisor, or a Hilbert--Pólya operator.  Hence
\[
 (A0,A1,A2,A3,A4)=(\mathrm{FAIL},\mathrm{WEAK},\mathrm{FAIL},
 \mathrm{FAIL},\mathrm{NATURAL\ QUANTIZATION}),
\]
overall \texttt{ROUTE\_A\_REJECTED}, Route B false, under
\texttt{NO\_BAD\_EULER\_OR\_ROOT\_NUMBER}.

\paragraph{Exact certificate.}
The package records 180 rational sentinels across both regimes and 1,260
zero coefficient residuals.  A producer-independent checker passes 4,268
assertions; a separate SymPy path proves 25 generic identities; replay is
byte exact; 20 repaired-hash semantic corruptions and one stale-hash
corruption are rejected.  These rows audit conventions and do not replace
the all-parameter proof above.

\ifcase\CRevisionRound
\paragraph{Revision-round focus.} Round 0 freezes the two Jacobi regimes and
their common $4K$ vector period.  It does not yet use the full sampled-time
owner boundary.
\or
\paragraph{Revision-round focus.} Round 1 adds all axial tangent rates, the
four explicit heteroclinic branches, and both KKS cap-action charts.
\or
\paragraph{Revision-round focus.} Round 2 makes the fixed-set statement
iterate-exact and closes the positive-dimensional zeta obstruction,
degenerate boundaries, evidence audit, and strict Route-A verdict.
\fi

\paragraph{Limitations.}
$G=0$ is the stationary origin.  Spherical and symmetric tops violate the
strict inertia gaps and have equilibrium manifolds requiring separate
elementary formulas.  We do not continue the regular formula through
$k=1$, count a fixed circle as finitely many points, claim arithmetic
semantics, or claim literature novelty or external review.

\section*{Declarations}
\paragraph{Data and code availability.} Exact evidence and deterministic
code accompany this manuscript.
\paragraph{Ethics.} No human, animal, clinical, personal, or sensitive data
are used; approval is not applicable.
\paragraph{Contributions.} This anonymous certificate records formal analysis,
software, validation, drafting, and revision. AI systems are not authors.
\paragraph{Funding and conflicts.} No external funding is reported; no
conflict is known.
\paragraph{AI-use disclosure.} An AI coding assistant supported drafting and
exact-code development; it was not an external reviewer or independent error
process.

\begin{thebibliography}{2}
\bibitem{Celledoni2008} E. Celledoni, F. Fassò, N. Säfström, and A. Zanna,
``The exact computation of the free rigid body motion and its use in
splitting methods,'' \emph{SIAM J. Sci. Comput.} 30 (2008), 2084--2112.
DOI: 10.1137/070704393.
\bibitem{Pina2015} E. G. Pina,
``Drawing the free rigid body dynamics according to Jacobi,''
arXiv:1505.06186 (2015).
\end{thebibliography}
\end{document}
